% Part: incompleteness
% Chapter: theories-computability
% Section: introduction

\documentclass[../../../include/open-logic-section]{subfiles}

\begin{document}

\olfileid{inc}{tcp}{int}

\olsection{પરિચય}

\begin{editorial}
આ વિભાગ ફરીથી લખવો જોઈએ.
\end{editorial}

આપણી પાસે નીચેની બાબતો છે:
\begin{enumerate}
\item $\Th{Q}$માં વિધેય નિરૂપ્ય
  હોવાનો અર્થ આપતી વ્યાખ્યા
  (\olref[req][int]{defn:representable-fn})
\item $\Th{Q}$માં સંબંધ નિરૂપ્ય
  હોવાનો અર્થ આપતી વ્યાખ્યા
  (\olref[req][rel]{defn:representing-relations})
\item $\Th{Q}$માં નિરૂપ્ય વિધેયો
  બરાબર સંગણનીય વિધેયો જ છે
  એવું કહેતો પ્રમેય
  (\olref[req][int]{thm:representable-iff-comp})
\item $\Th{Q}$માં નિરૂપ્ય સંબંધો
  બરાબર સંગણનીય સંબંધો જ છે
  એવું કહેતો પ્રમેય
  (\olref[req][rel]{thm:representing-rels})
  \sourcecorrection{OLINC-038}{સ્થિર અંગ્રેજી
  મૂળની આ છેલ્લી યાદી-બાબતમાં
  સંદર્ભની શરૂઆતનો ગોળ કૌંસ
  ગાયબ છે; અગાઉની ત્રણ
  બાબતોની જેમ અહીં બંને
  કૌંસ આપ્યા છે.}
\end{enumerate}

\emph{સિદ્ધાંત} એ નિગમન હેઠળ
બંધ વાક્યોનો ગણ છે; એટલે
$T$માંથી $!A$ સાબિત થાય
ત્યારે $!A$ એ~$T$માં જ
હોય છે. સિદ્ધાંતને વાક્યોનો
એવો સંગ્રહ સમજવો સારું છે,
જેમાં તે વાક્યોમાંથી નીકળતાં
બધાં તારણો પણ સમાવિષ્ટ છે.
હવેથી $\Th{Q}$ વડે
\olref[req][int]{sec}ની
આઠ સ્વયંસિદ્ધિઓમાંથી નિગમ્ય
વાક્યોના ગણરૂપ~\emph{સિદ્ધાંત}
દર્શાવીશું. યાદ રાખો કે
$\Th{Q}$ની ભાષાનાં સૂત્રોને
સંખ્યાઓ વડે સંકેતિત કરી
શકાય છે; $!A$ એવું સૂત્ર
હોય, તો $\Gn{!A}$ એ~$!A$ને
સંકેતિત કરતી સંખ્યા છે.
આ સંકેતીકરણને ધ્યાનમાં રાખીને
હવે વિવિધ સૂત્રગણો સંગણનીય
છે કે નહીં તે પૂછી શકીએ.

\end{document}
